\chapter*{Introduction}
\addcontentsline{toc}{chapter}{Introduction}
\markboth{Introduction}{Introduction}

\label{chap:intro}
\setlength{\parskip}{1em}

Modern cyber conflicts have escalated into a condition of constant cyberwar, driven by countless threats. Unlike traditional military conflicts, this warfare is taking place within a digital space and often goes unnoticed by the general public due to their lack of awareness and competence \cite{almeida2025comparative}. The current policy of major states is dictated by the principle of "peace through strength", prioritizing the rule of power over ethical considerations and international law. It leads to an intensified arms race between multiple actors, such as government agencies, armies, corporations, and private groups. Under these circumstances of ongoing confrontations between defenders and attackers, ensuring security is only a continuous process rather than an absolute guarantee \cite[pp.~4--6]{safitra2023counterattacking}. So, the industry of information technology (IT) should be engaged in systematic risk assessment and apply reactive patches as problems emerge, because the number of potential threats and vulnerabilities is inherently unknown and cannot be fully identified or eliminated upfront. Moreover, the increasing rate of automation and digitalization exposes associated risks for critical infrastructure, where their compromise could create dangerous situations and collateral damage for citizens \cite[pp. 1–2]{nist2018framework}. For example, CISA (the Cybersecurity and Infrastructure Security Agency of the USA) adopted a classification that identifies 16 critical infrastructure sectors, including the context of IT infrastructures \cite[pp. 6–7]{cisa2025infrastructure}. This sector is composed of corresponding services and assets (e.g., maintenance personnel, electricity, networks, equipment, hardware, software), providing an operational basis or computational environment to utilize for information systems (IS) \cite{qatawneh2024mediating}.

Fundamentally, attacks targeting critical IT infrastructure are driven by diverse motives and reasons, overlapping with other industries \cite{lehto2022cyber}. Historical events have already shown that conflicts of national interests might trigger a security breach on highly protected facilities (e.g., nuclear power plants) \cite{makrakis2021industrial}. For instance, Riggs et al. examined statistics on global cyberattacks against critical infrastructure from 2006 to 2022 and developed a prediction model for the following 5 years that shows an exponential growth of incidents, particularly within the military sector \cite{riggs2023impact}. On the other hand, the race of AI and associated information technologies demand a rapid expansion of computational resources in IT infrastructures, which increases the demand on data centers, stimulating competitors and other third parties to compromise or abuse that computational power for their own benefit \cite{yigit2024critical}. The technofeudalism system of major corporations dictates conditions in which their cloud services eliminate individual ownership through vendor lock-in ecosystems and subscription-based models. These market conditions drive manufacturers and suppliers to operate as cartels, dictating prices and monopolizing across various enterprise segments. For example, Pilz et al. highlight a trend based on data from 2019 to 2025 years, showing that the computational performance of leading AI supercomputers doubles every 9 months, whereas their power requirements double every 1 year \cite[Sec. 3.2-3.3]{pilz2025trends}. Moreover, the total electricity capacity required even by a single leading AI data center may reach 9 GW (gigawatts) by 2030 in projections, given the continuing pace of growth \cite[Sec. 4.1-4.2]{pilz2025trends}. Collectively, these factors provoke market manipulations and stimulate growth of prices with shortages of critical resources (e.g., energy and raw materials, hardware components, qualified personnel, water for cooling, etc.). The over-reliance on centralized data centers represents single points of failure that are vulnerable to physical damage or service disruptions caused by state policies. This is particularly critical given that statistics of 2025 show 75\% of data centers are concentrated within the United States of America (USA), with the remainder distributed in other countries \cite[Sec. 3.6]{pilz2025trends}. Also, many hardware technologies and their production remain strictly proprietary, resulting in a lack of competition between actors and reducing overall production rates. Specifically, the USA's national policy has already undergone three phases since 2018 for partial export bans on critical technologies targeting China and other market participants \cite[Ch. 3]{yu2025blockade}. Additionally, the increasing interconnectivity of IS, especially through Internet of Things (IoT) devices, causes fragility by exposing multiple points of failure and allows malicious actors or even device manufacturers to gain remote access to them \cite{vardakis2024review,rauhala2022physical}. The industry faces various issues, including state-mandated backdoors, supply chain bottlenecks, planned obsolescence, and the accumulation of electronic waste.

For these reasons, independent nations should pursue technological autonomy by developing local production in software and hardware without becoming hyper-dependent on untrusted partners, while still leveraging global experience and standards. Especially, military contexts of combat operations require portable and autonomous systems to mitigate the risk of lost or disrupted communications. Historically, this necessity drove the development of networks such as ARPANET and Usenet, born from the synergy of academic research and government-military collaboration \cite{leiner2009brief}. Therefore, this study attempts to construct security measures for the IT infrastructure, emphasizing the process of collecting and accounting of its internal resources or components for reliability and early insider threat prevention. Also, it focuses on investigating how small-scale groups can compete in the AI sphere, combining cost-effective or low-power devices with data privacy without exposure to third parties. Firstly, the literature review covers the historical development of IT infrastructure, complexity management, issues with software composition, security aspects, application of decentralized networks for general-purpose tasks, and democratization of AI. Secondly, the methodology specifies the hypotheses and provides a framework for conducting customizable analyses within sandboxed environments, and shows the concept of the proposed distributed network by covering its tokenomics, resistance, and fairness to prevent centralization of nodes in pools. Thirdly, the results part examines features of the proposed solution and describe our dataset of Android-compatible models of smartphones with their characteristics. Finally, the conclusion summarizes findings and potential drafts for future works.
